\documentclass{article}
\usepackage{amsmath,url}
\title{Generating a blocked source from motion, a stop, and live demand}
\begin{document}
\maketitle
\paragraph{Scope.}
This is a small mechanistic extension of the existing pinned MRWS development experiment, not a complete fault theory. The primitive fault opportunities are a battery stop and delayed or absent disclosure of that same stop. The simulator generates actual locations, transport, shelf identities, and order progress. A blocked shelf is not another free fault switch. The route and commissioning mechanisms still determine which events are possible. A repair changes which histories can happen even though the two fault opportunities are fixed.

\paragraph{Component interfaces.}
The native plant executes the motion and stock updates. A location monitor folds create and move events by actor, and records whether the latest motion was protected by the guard. A stock monitor folds physical shelves, cargo gains and losses, and the goal's delivered prefix. A fault monitor freezes the stopped actor and associates a possibly absent notice with that actor's actual stop and time. At a shared stop event the monitors join on actor, event order and cut. The derived hazard is occupation of the sole shelf supplying a still-unsettled requested item. It is not sufficient for service loss: 64 stop records have such an occupation while another surviving carrier already physically holds that item. Neither a physical copy nor observation of occupation grants the right to use the copy. A separate live commissioning fact comes from the scheduler's actual order membership, not from stock or geometry. The monitors are analysis tools: their joined global state is not exposed to a live controller that only receives notices.

The audited join is checked against independently read simulator position, inventory, and goal prefix at the actual native fault. A caught instrumentation error illustrates the point: testing membership in the residual of \emph{all} shelf labels had classified a stop on item 0 as blocking the order. Item 0 exists physically but is not requested in this nine-item task. The revised condition matches the actual request. Ignoring event identity, using only fault notices, or letting shelf obstruction be an independent switch cannot validate this interpretation. The exact programming interfaces and checks are in the replay. This is trace generation by a pinned native plant plus interpretation, not a replacement semantics for all MRWS operations.

\paragraph{A generated exclusion.}
Fix the same order instance, a monotonically growing delivered prefix, a stopped body that retains its pose, and the case that the last actual move before the stop used the home-route guard. Its destination was outside the cells needed at the movement cut. Since the set of needed cells only shrinks before the stop, that last move cannot generate this blocked-source event. This proof does not apply to a stop before the first protected move, after a pick or other unprotected move, or after changing the order. It neither removes the stop opportunity nor restores any missing administrative power. Across the wide and scoped variants, all 236 stops after a protected last move satisfy that exclusion. Earlier stops still occupy required shelves in some histories. Lean checks the logical exclusion and two observation/identity consequences, not the native runtime.

\paragraph{The same repair can create a different shortage.}
The source generates the first stop from the same reachable initial cut and admits a possible second stop one through eight ticks later, then all specified notice delays. The old audit used only lags 1 and 4. Compare the original planner, a route guard on every home move, and a later predeclared test that restricts the guard to carriers no longer commissioned for the order. Among 264 paired opportunity settings, the broad guard restores service in 33 runs that originally did not finish, but loses three previously successful runs. In the latter runs no second-stop shelf obstruction is present in either arm. The same second stop and delivered notices lead to a different state at the replan. In the broad arm, the retained-load rule refuses because none of the live carriers is empty: one already holds items 3, 2, and 1, and the other already holds item 9; nothing has yet been delivered. The narrow guard restores all three (finish at tick 77) and retains 24 gains over the original, but gives up nine of the broad guard's newly successful cases. No configuration is selected after seeing individual outcomes, and neither guard dominates the other. The changed action altered availability of a spare rather than the fault rate. The stopped positions and scores are in the archive.

\paragraph{Calibration limits and comparator.}
The native replays cover 792 histories and 113,009 monitored events. The broad guard changes the physical source blockage in 60 matched cases with the same notice sequences, including 52 outside the earlier pair of lag values. These are counts over authored inputs, not rates. The causal comparison varies actual allowed route changes and checks source state; it does not infer mechanisms from a table of labeled faults. A direct complete native trace with all the same facts can derive the same hazard. A general compositional fault generator would also need an explicit rule for composing new mechanisms and validating their admissible exogenous inputs. This study does not supply that rule, independent external validation, or contract acceptance.
\end{document}
